% Responsable: Jerónimo
% Fuentes: m34_clusters.py, CLAUDE.md de Modelo3 §8 y §7.4, m34_placebo.py,
% Seleccion_Variables_Restricciones_Modelo4.pdf
% Tesis: si el Estado conoce el TIPO de hogar (no el hogar), ¿cuánto mejora? Y, con el
% placebo: ¿la mejora es por la información demográfica o solo por partir más fino?
\subsection{Modelo 4: focalización mixta geográfica y demográfica}\label{sec:m4}

El Modelo 4 agrega al Modelo 3 una segunda fuente de información que el Estado sí
tiene: las características demográficas de los hogares, del tipo que registra el
SISBÉN (tamaño, edades, jefatura). Sigue sin saber quién está privado en qué, pero
ahora puede dirigir algunos programas a \emph{tipos} de hogar y no solo a zonas.

\subsubsection{Bienes públicos y bienes privados}

Siguiendo a \citet{hamie2024}, los indicadores se dividen en dos grupos. Los de
\emph{bien público}, $\Vpub$, se proveen con infraestructura que sirve a todos los
hogares de un territorio y por eso solo pueden focalizarse geográficamente. Los de
\emph{bien privado}, $\Upriv$, se entregan hogar por hogar y pueden dirigirse a un
tipo de hogar. Tomamos
\[
\Vpub=\{\text{acueducto},\ \text{alcantarillado}\},\qquad
\Upriv=\Ind\setminus\Vpub\ \ (\text{los otros 13 indicadores}).
\]
La partición se apoya en los datos. La Tabla~\ref{tab:discrimina} compara qué
tanto varía cada privación entre zonas y entre tipos de hogar, medido como la
razón entre la incidencia en la celda más privada y en la menos privada. El
acueducto y el alcantarillado varían más entre zonas que entre tipos, así que
la geografía es la forma natural de focalizarlos. La inasistencia escolar, el
hacinamiento y la atención integral a la primera infancia varían mucho más entre
tipos de hogar; la última solo existe en hogares con menores de cinco años. En
las privaciones de salud, que son las primeras que atienden los modelos
(Sección~\ref{sec:res34}), las dos particiones discriminan casi lo mismo.

\begin{table}[ht]
\centering\footnotesize
\caption{Razón entre la incidencia de la privación en la celda más privada y en
la menos privada, por zona y por tipo de hogar ($k=4$; personas expandidas).}
\label{tab:discrimina}
\begin{tabular}{@{}lrr@{}}
\toprule
Privación & Entre zonas & Entre tipos de hogar \\
\midrule
Acueducto (bien público)        & \cDiscAcuZona{} & \cDiscAcuClu{} \\
Alcantarillado (bien público)   & 4,7  & 3,2 \\
Inasistencia escolar            & 4,8  & \cDiscInasClu{} \\
Hacinamiento                    & 2,7  & 26,6 \\
Atención integral (0--5 años)   & 2,2  & solo en hogares con niños \\
Aseguramiento en salud          & 1,9  & 1,9 \\
Barreras de acceso a salud      & 2,7  & 1,8 \\
\bottomrule
\end{tabular}
\end{table}

\subsubsection{Tipos de hogar}

Los tipos $\Tip$ se construyen con $k$-medias sobre ocho rasgos que el Estado
observa sin necesidad de una encuesta de privaciones:
\begin{itemize}
\item número de personas;
\item número de menores de 5 años, de personas entre 5 y 17 años y de mayores de 65;
\item tasa de dependencia;
\item edad, sexo y nivel educativo del jefe del hogar.
\end{itemize}
Los rasgos se estandarizan, cada hogar se pondera por su factor de expansión y la
semilla queda fija para que el resultado sea reproducible.

Hay dos decisiones deliberadas. Primero, \textbf{ninguna privación del IPM entra
en el \emph{clustering}}: si los tipos se formaran con las privaciones, el
planificador estaría usando justo la información que el modelo supone que no
tiene, y el Modelo 4 dejaría de simular información imperfecta. Segundo, la base
del IPM no trae ingreso ni gasto, así que el \emph{proxy} de ingreso que usa el
paper no está disponible. Ambas diferencias con el paper están en el
Apéndice~\ref{app:desv}.

Con $k=4$ los tipos tienen una lectura clara (Tabla~\ref{tab:clusters}) y
separan la pobreza sin haberla usado: la incidencia del tipo más pobre es
\cSeparaPobreza{} veces la del menos pobre. Parte de esa separación es
estructural, y lo decimos: la privación de atención integral solo puede ocurrir en
hogares con niños pequeños.

\begin{table}[ht]
\centering\small
\caption{Tipos de hogar del Modelo 4 ($k=4$).}\label{tab:clusters}
\begin{tabular}{@{}lrr@{}}
\toprule
Tipo & Hogares en la muestra & Incidencia $H$ \\
\midrule
Adultos sin dependientes & 1.389 & 3,7\,\% \\
Adultos mayores          & 1.038 & 11,8\,\% \\
Con hijos escolares      &   849 & 9,1\,\% \\
Primera infancia         &   456 & 16,1\,\% \\
\bottomrule
\end{tabular}
\end{table}

\paragraph{Elección de $k$.}
El coeficiente de silueta casi no cambia entre $k=2$ y $k=8$ (de \cSiluetaMin{}
a \cSiluetaMax{}), así que no sirve para elegir. Elegir $k$ es, en el fondo, una
decisión sobre qué tan fino es el sistema de información del Estado, y un $k$
muy grande termina simulando información individual: con $k=8$ hay celdas de
menos de 10 hogares privados, donde el sorteo deja de representar una
focalización por tipo. Adoptamos dos criterios:
\begin{enumerate}
\item que menos del 15\,\% de los pares (indicador, celda) tenga menos de 10
  hogares privados en la muestra;
\item que al pasar a $k+1$ la ganancia sea marginal.
\end{enumerate}
$k=4$ cumple los dos y $k=5$ queda como alternativa. El análisis también muestra
que no cualquier partición demográfica sirve: con $k=2$ los tipos se separan por
tamaño del hogar, que casi no se relaciona con la pobreza, y el Modelo 4 resulta
\emph{peor} que el Modelo 3 (Sección~\ref{sec:res34}).

\subsubsection{Formulación}

El Modelo 4 es el Modelo 3 con otra regla para asignar cada par a su celda de
focalización:
\[
g(i,j)=\begin{cases} d_i & j\in\Vpub,\\ t_i & j\in\Upriv,\end{cases}
\]
de modo que el gasto, la cobertura y el sorteo
\eqref{eq:m3-gasto}--\eqref{eq:m3-sorteo} se definen por par $(j,g)$, y el
objetivo es
\begin{equation}\label{eq:m4-obj}
\min\ \sum_{j\in\Upriv}\sum_{t\in\Tip}\Esf_{jt}+\sum_{j\in\Vpub}\sum_{d\in\Zon}\Esf_{jd},
\end{equation}
que corresponde a OBJ~4 de \citet{hamie2024}. Las restricciones de pobreza y la
meta \eqref{eq:m3-r4}--\eqref{eq:m3-meta} no cambian, porque el puntaje $s_i$
depende de todos los indicadores sin importar cómo se focalizó cada uno. La
propiedad de prefijo de la Proposición~\ref{prop:cortes} solo usa que la
cobertura es común dentro de una celda, así que vale igual dentro de un tipo de
hogar, y la reformulación por cortes se aplica sin cambios. El Modelo 4 tiene
\cParesMcuatro{} pares y entre \cCortesMcuatro{} cortes admisibles según la
réplica, y cada instancia se
resuelve en segundos.

\subsubsection{Placebo: ¿información demográfica o partición más fina?}

El Modelo 4 difiere del Modelo 3 en dos cosas a la vez. Usa información
demográfica, pero también reparte los 13 indicadores privados en 4 celdas en vez
de 3, y una partición más fina focaliza mejor aunque no tenga ningún contenido. El
caso límite lo muestra: con un hogar por celda, el Modelo 3 coincide con la
información perfecta (Sección~\ref{sec:diseno}). Para separar los dos efectos
construimos un \emph{placebo}: un Modelo 4 en el que las etiquetas de tipo se
permutan al azar entre los hogares. Así se conservan exactamente el número y el
tamaño de las celdas, pero desaparece cualquier relación entre el tipo y las
privaciones. Cada réplica usa su propia permutación y el mismo sorteo $R$ que el
Modelo 3 y el Modelo 4, de modo que la comparación es pareada. Si el Modelo 4
no supera al placebo, su ventaja sobre el Modelo 3 se debe a la granularidad y no
a la demografía.
